\begin{textsample}{POS Dim 4 – vem\_ed\_15 – Score 3.00 – vem\_ed\_15/t067.txt}  \label{ex:f4_pos_118}
EXPERIÊNCIAS DE DOCENTES DAS FATECS Professores de Fatecs compartilharam suas práticas em PCI: Rose Maria Duda (Fatec Jaboticabal) falou sobre colaborações com a Universidad El Bosque e o Instituto Superior de Educación Rural (ISER), ambos na Colômbia.
Alunos de Gestão Ambiental da Fatec Jaboticabal e de Engenharia Ambiental da Universidad El Bosque detectaram problemas e buscaram soluções envolvendo temas como qualidade da \textbf{água}, poluentes na \textbf{água} e no esgoto e \textbf{rios} fronteiriços entre os dois países.
A outra colaboração envolveu estudantes de Tecnologia em Biocombustíveis (Fatec Jaboticabal) e de Tecnologia em Processos Agroindustrial (ISER) e contou com o apoio da profa.
Camila Carla Guimarães.
As equipes mistas estudaram possibilidades de utilização de resíduos na produção de óleo de palma e de açúcar e etanol da cana-de-açúcar.
Leysiane Cristina Pavani Pazini (Fatec Garça) contou sobre o projeto com Setsunan University, no Japão.
Esse PCI voltado ao estudo dos Objetivos de Desenvolvimento Sustentável (ODS), com a prática da língua inglesa por parte de falantes não-nativos, teve o diferencial de ser realizado exclusivamente com encontros síncronos.
O fuso de 12 horas, no caso, foi um aliado, porque Leysiane leciona à noite e Curtis, de manhã.
Jorge Fernando Tenório e Mauro Yuji Ohara (Fatec Franco da Rocha) relataram o PCI realizado com Malgorzata Marchewka, professora da Universidade de Economia da Cracóvia, na Polônia.
Jorge é professor de inglês e Mauro leciona gestão de projetos.
Os estudantes desenvolveram competências interculturais, linguísticas, digitais e de trabalho em equipe, aplicando na prática conceitos relativos ao gerenciamento de projetos: distribuição das tarefas, organização das etapas e \textbf{prazos} para entrega das atividades.
Os grupos de alunos abordaram problemas reais vividos por eles.
Uma das equipes surpreendeu, realizando um show ao vivo no Instagram para angariar doações para as crianças vítimas da guerra na Ucrânia.
Mais sobre a iniciativa no site Cesu: https://cesu.cps.sp.gov.br/pci-entre-fatec-franco-da-rocha-e-uek-polonia-gera-live-para-angariar-doacoes-para-criancas-e-jovens-da-ucrania/.

% matched lemmas: prazo, rio, água
\end{textsample}
